\begin{lemma}[Weak closedness in $L^2(\gamma_n)$ of the nonnegativity constraint]
Let $(v_k)_{k \geq 1} \subseteq L^2(\gamma_n)$ satisfy $v_k \geq 0$
$\gamma_n$-almost everywhere for every $k$, and suppose $v_k \rightharpoonup v$
weakly in $L^2(\gamma_n)$, i.e.
$\langle v_k, \varphi\rangle \to \langle v, \varphi\rangle$ for every
$\varphi \in L^2(\gamma_n)$.  Then $v \geq 0$ $\gamma_n$-almost everywhere.
Consequently the admissible cone $\mathcal{K}$ of \noderef{ObstacleCone} is
closed under weak $L^2(\gamma_n)$ limits of sequences with uniformly bounded
energy, and no pointwise convergence is needed.
\end{lemma}
\begin{proof}
Let $\varphi \in L^2(\gamma_n)$ with $\varphi \geq 0$ $\gamma_n$-a.e.  By the
implication $(1) \Rightarrow (2)$ of \noderef{NonnegTestCharacterization}, applied
to each $v_k$, we have
\[
  \langle v_k, \varphi\rangle = \int v_k \varphi \, d\gamma_n \geq 0
  \qquad\text{for every } k .
\]
By the weak convergence hypothesis, $\langle v_k, \varphi\rangle \to \langle v, \varphi\rangle$.
A limit of a sequence of nonnegative reals is nonnegative, so
$\int v \varphi \, d\gamma_n = \langle v, \varphi\rangle \geq 0$.

Since $\varphi$ was an arbitrary nonnegative element of $L^2(\gamma_n)$, the
implication $(2) \Rightarrow (1)$ of \noderef{NonnegTestCharacterization} applies
to $v$ and yields $v \geq 0$ $\gamma_n$-almost everywhere.
\end{proof}
